%% Chapter 6 — Real-Time Feedback and Mid-Circuit Measurement
\section{Real-Time Feedback and Mid-Circuit Measurement}
\label{sec:feedback}

A key motivation for FPGA-based control is the ability to execute
\emph{real-time} conditional logic: perform a measurement, make a decision, and
apply a corrective pulse, all on the FPGA.
This section describes the tProcessor instruction set, demonstrates the
conditional branching primitive, and outlines the active qubit reset protocol.

\subsection{tProcessor Architecture}
\label{subsec:tproc_arch}

The tProcessor is a soft-processor synthesised in the FPGA fabric, with a
custom instruction set tailored to pulse generation and readout sequencing.  In
the original (v1) design~\cite{Stefanazzi2022}, on whose demo firmware the
branching demonstration below was run, each
channel page holds 32 general-purpose 64-bit registers, with cross-page registers
for inter-channel synchronisation.  Dedicated instructions fire the DAC generators
or arm the ADC acquisition windows, while timing instructions advance the schedule:
\texttt{synci} moves the tProcessor clock without touching the channel timelines,
\texttt{sync\_all} aligns all channels, and \texttt{wait\_all} stalls until the
readout windows close.  The primitive that matters for feedback, common to v1 and
the tProcessor v2 build used for the readout experiments
(\texttt{qick\_processor} rev 21: \SI{200}{\mega\hertz} core execution clock,
pulse timing resolved on a \SI{614.4}{\mega\hertz} timing clock), is the conditional
branch \texttt{condj page, reg\_a, op, reg\_b, label}, which jumps to \texttt{label}
when \texttt{reg\_a op reg\_b} holds.

All branching targets are resolved at \emph{compile time} by the Python
assembler: \texttt{label()} records instruction positions as named addresses and
jump targets are backpatched before the program is loaded, so at run time only
integer addresses exist.

\subsection{Conditional Logic Demonstration}
\label{subsec:cond_logic}

The \QICK{} Demo 3 notebook demonstrates the core conditional-branch primitive:
\begin{minted}[fontsize=\small,linenos]{python}
def body(self):
    cfg = self.cfg
    self.trigger(adcs=[cfg["ro_ch"]],
                 adc_trig_offset=cfg["adc_trig_offset"])
    # condj: jump to SKIP_PULSE if number < threshold
    self.condj(0, 2, '<', 1, 'SKIP_PULSE')
    self.pulse(ch=cfg["res_ch"])          # skipped if condition true
    self.label('SKIP_PULSE')
    self.wait_all()
    self.sync_all(self.us2cycles(cfg["relax_delay"]))
\end{minted}

The demo uses the v1-style programming interface (\texttt{body}, \texttt{synci};
Section~\ref{subsec:tproc_arch}), not the tProcessor v2 API used elsewhere
(Section~\ref{subsec:qick_stack}).  Registers 1 and 2 are pre-loaded with the threshold and the test value via
\texttt{regwi}.  When \texttt{number=100, threshold=50}: the condition
\texttt{100 < 50} is false, the pulse fires, and the ADC captures a loopback
signal.  When \texttt{number=10, threshold=50}: the condition is true, the pulse
is skipped, and the ADC trace is flat noise.

\paragraph{Timing subtlety.}
Inside a \texttt{body()} with an open ADC acquisition window, time must be
advanced with \texttt{synci}, not \texttt{sync\_all}: \texttt{sync\_all} advances
all channel timelines to the furthest point, the end of the ADC window, so
subsequent pulses would land after it closes; \texttt{synci} advances only the
tProcessor's internal counter.

\subsection{\texorpdfstring{$\pi$}{π} Pulse Calibration}
\label{subsec:pi_pulse}

A $\pi$ pulse rotates the qubit state by 180° on the Bloch sphere:
$\ket{0} \to \ket{1}$ and $\ket{1} \to \ket{0}$.  It is a fundamental primitive
for qubit initialisation (active reset) and single-qubit gates.  Calibrating it
requires a connected qubit chip; this section describes the procedure that would
be followed once the ZCU216 is connected to a device.

The $\pi$ pulse amplitude is found via a \emph{Rabi experiment}: the qubit drive
pulse amplitude is swept from zero, and the readout signal oscillates between the
$\ket{0}$ and $\ket{1}$ IQ clusters as the qubit is driven through successive
half-turns on the Bloch sphere.  The first amplitude at which the signal fully
crosses to the $\ket{1}$ cluster is the $\pi$ amplitude.  The Rabi sequence is:
\begin{enumerate}
  \item Prepare qubit in $\ket{0}$ (wait $\gg T_1$).
  \item Apply qubit drive pulse at frequency $\omega_{01}$ with variable amplitude
        $A$ and fixed duration.
  \item Read out via the resonator.
\end{enumerate}

The qubit drive tone is at $\omega_{01}$, distinct from the resonator frequency
$\omega_r$ used for readout.  For transmons $\omega_{01}$ is typically
\SIrange{4}{8}{\giga\hertz}; for fluxonium biased near half flux it is far lower,
of order \SIrange{0.1}{1}{\giga\hertz}, which changes the drive-line requirements
but not the sequence described here.  Two separate generator channels
are therefore required: one for readout, one for qubit control.  The
\QICK{} firmware supports this natively: the tProcessor can address both channels
independently with nanosecond timing precision.

\subsection{Active Qubit Reset}
\label{subsec:active_reset}

Active reset rapidly returns the qubit to $\ket{0}$ by measuring its state and
applying a corrective $\pi$ pulse only if the qubit is found in $\ket{1}$~\cite{Reed2010}.  The
\QICK{} tProcessor v2 has all the primitives required to implement this entirely
on-FPGA: the conditional branch demonstrated in Section~\ref{subsec:cond_logic}
provides the decision logic, and the two-channel architecture supports
simultaneous readout and drive.  The following protocol, assembled from these
primitives, would be deployed once the board is connected to a qubit chip:
\begin{enumerate}
  \item \textbf{Dispersive readout}: fire the readout tone and acquire the
        accumulated IQ result.
  \item \textbf{Thresholding}: the readout firmware compares the IQ result against
        a pre-calibrated threshold and writes the binary decision into a
        tProcessor register.
  \item \textbf{Conditional $\pi$ pulse}: \texttt{condj} checks the decision
        register and applies a calibrated $\pi$ pulse on the qubit drive channel
        only if the qubit is in $\ket{1}$.
  \item \textbf{Wait}: allow any residual photons in the resonator to decay before
        starting the experiment.
\end{enumerate}

Because the entire decision-to-pulse path would run inside the tProcessor, the
latency between completing the readout window and firing the corrective pulse is
determined solely by FPGA clock cycles.  For the tProcessor v2 build used here,
with its \SI{200}{\mega\hertz} core execution clock, this latency is of order tens
of nanoseconds.

\subsection{Latency Budget}

Table~\ref{tab:latency} gives an estimated latency breakdown for the active reset
sequence.

\begin{table}[ht]
\centering
\caption{Estimated latency contributions in the active reset loop.}
\label{tab:latency}
\begin{tabular}{lrl}
\toprule
Stage & Latency & Notes \\
\midrule
Readout pulse duration     & \SIrange{0.5}{2}{\micro\second}   & Depends on $\kappa$ \\
ADC integration window     & \SIrange{0.5}{2}{\micro\second}   & \\
FPGA decision (threshold)  & $\sim$\SI{10}{\nano\second}        & A few core cycles \\
Conditional branch         & $\sim$\SI{5}{\nano\second}         & 1 core cycle (\SI{200}{MHz}) \\
$\pi$ pulse duration       & \SIrange{50}{500}{\nano\second}    & Set by drive strength; $\ll T_1$ \\
\midrule
Total (excl. readout)      & $<$\SI{1}{\micro\second}           & \\
\bottomrule
\end{tabular}
\end{table}

The dominant latency is the readout integration window; the digital decision and
conditional fire together contribute less than \SI{100}{\nano\second}, well within
qubit coherence times of \SIrange{10}{100}{\micro\second} for state-of-the-art
fluxonium devices.
